% Summarize the key findings, contributions, and implications of the project.

This paper presented \textit{Algiers in 24 Hours}, a tourist route
optimisation system grounded in the Time-Constrained Orienteering
Problem with Time Windows. Six algorithmic approaches spanning the
full spectrum from deterministic construction to exact
branch-and-bound were implemented, integrated into a common model
layer, and evaluated on a curated dataset of 56 Algiers landmarks
with verified GPS coordinates, opening hours, and interest scores.

The principal finding is that Tabu Search with Strategic
Oscillation achieves solution quality equivalent to GRASP (98.7
interest score, 11 landmarks) at a fraction of the computational
cost (0.37 versus 10.18 seconds), establishing it as the
recommended solver for real-time deployment contexts. Simulated
Annealing with the Cauchy acceptance criterion provides a
competitive alternative (96.9) with moderate runtime. The
underperformance of CPLEX relative to leading metaheuristics
underscores the known intractability of the OPTW for branch-and-bound
approaches under realistic time budgets, and motivates the design
of hybrid exact-heuristic methods as a direction for future
investigation.

Beyond algorithmic contribution, this work produces a reusable,
open-source problem model class for the Algerian tourist
orienteering context --- a domain absent from existing benchmark
libraries --- and a full-stack web application through which
algorithm behaviour can be explored interactively by non-expert
users. The gap between the best heuristic score and the CPLEX
incumbent across all experiments defines a concrete quality target
for future work. Closing that gap through Adaptive Large
Neighbourhood Search or learned construction heuristics, while
extending the formulation to multi-day and stochastic settings,
constitutes the principal agenda for subsequent research.
